\problem{Q13}{感知机：按样本顺序学习 OR}{正面第2栏／P01；可分性见 P02}
\qtext{输入按 $(0,0),(0,1),(1,0),(1,1)$ 顺序反复出现，目标依次为 $z=0,1,1,1$。规定 $u=w_1x_1+w_2x_2+b$，$\hat z=H(u)$，其中 $H(u)=1$ 当 $u\ge0$，否则为0。初值 $w_1=w_2=b=0$，学习率 $\eta=1$。用 $\tilde w\leftarrow\tilde w+\eta(z-\hat z)\tilde x$ 更新，$\tilde x=(x_1,x_2,1)\T$。完成第一轮，判断是否收敛，再继续到一整轮没有误分类。}
\route{圈出“\kw{阈值输出、0/1目标、逐样本更新、顺序}”，查 P01。这里用类别差 $z-\hat z$；不能用 ADALINE 的连续残差 $z-u$。P02 提醒：OR 可分，标准感知机训练有有限步收敛保证。}
\given{$x\in\R^2$，$w\in\R^2$，$b\in\R$；增广后 $\tilde x,\tilde w\in\R^3$。阈值形式为 $w\T x-\theta$ 时，$b=-\theta$。每个样本先用当前全部旧参数预测，再更新全部分量。}
\solve
第一轮的每一行都从上一行更新后的参数开始：
\[
\begin{array}{c|r|c|r|c}
 x&u&\hat z&z-\hat z&\text{更新后 }(w_1,w_2,b)\\\hline
 (0,0)&0&1&-1&(0,0,-1)\\
 (0,1)&-1&0&1&(0,1,0)\\
 (1,0)&0&1&0&(0,1,0)\\
 (1,1)&1&1&0&(0,1,0)
\end{array}
\]
例如第二行的代入是
\[
 (0,0,-1)+(1-0)(0,1,1)=(0,1,0).
\]
第一轮结束后，$(0,0)$ 仍给 $u=0$、预测1，因此尚未收敛。继续按同一顺序扫描：
\begin{itemize}
\item 第二轮遇 $(0,0)$ 时更新为 $(0,1,-1)$；遇 $(0,1)$ 无需更新；遇 $(1,0)$ 时 $u=-1$、误分类，更新为 $(1,1,0)$；最后一点正确。
\item 第三轮遇 $(0,0)$ 时更新为 $(1,1,-1)$；其余三点都正确。
\item 第四轮四点全部正确，没有任何更新，才满足停止检查。
\end{itemize}
\ans{第一轮参数为 $(0,1,0)$，尚未收敛；最终可得 $(1,1,-1)$，判别式 $\hat z=H(x_1+x_2-1)$，共发生5次非零更新。}
\qcheck{最终四点分数为 $-1,0,0,1$，按 $H(0)=1$ 得 $0,1,1,1$。更新次数不等于样本访问次数；最后还需扫描一轮确认。}
\pitfall{把 $H(0)$ 当0会改变整条轨迹。漏掉增广输入的末项1，就无法更新 bias；若直接更新阈值，应为 $\theta\leftarrow\theta-\eta(z-\hat z)$。}
\transfer{换样本顺序或初值，路径和最终分隔线可能改变；仍先预测再纠错并检查整轮。若改成 XOR，先查 P02 判断结构是否有解，增加轮次不改变单阈值的表示限制。}

\problem{Q14}{为什么 XOR 不能由一个感知机实现？}{正面第2栏／P02；多层构造见正面第3栏／B01}
\qtext{令 XOR 目标为 $z(0,0)=z(1,1)=0$，$z(0,1)=z(1,0)=1$。单个感知机使用 $\hat z=H(w_1x_1+w_2x_2+b)$，$H(0)=1$。用不等式证明不存在正确分类四点的参数。再解释：“感知机不收敛只是学习率太小”“有限数据可分就能严格分开”“收敛后一定得到最大间隔”三句话。}
\route{圈出“\kw{XOR、single perceptron、linear separability、convergence}”，查 P02。题目先问模型能否表示；不应先套 LMS 或 BP 调学习率。多层网络的非线性中间表示才可能改变可分性。}
\given{$u=w\T x+b$ 是仿射分数；$z=1$ 要求 $u\ge0$，$z=0$ 要求 $u<0$。参数 $w_1,w_2,b$ 是未知量，四个样本给四条约束。}
\solve
逐点写出条件：
\[
\begin{array}{c|c}
 (0,0):z=0&b<0\\
 (0,1):z=1&w_2+b\ge0\\
 (1,0):z=1&w_1+b\ge0\\
 (1,1):z=0&w_1+w_2+b<0
\end{array}
\]
两个正类条件相加，得 $w_1+w_2+2b\ge0$；两个负类条件相加，却得 $w_1+w_2+2b<0$，矛盾。因此参数不存在。

几何上，两个正类顶点连成的线段和两个负类顶点连成的线段共用中点 $(1/2,1/2)$。仿射分数在一条线段中点等于两端分数的平均；同一中点不可能同时属于正、负两侧。

三句判断分别为：
\begin{enumerate}
\item 第一项错误。对单阈值 XOR，失败来自表示能力；调整步长、训练更久不能产生不存在的解。一般数据不收敛也可能有实现或优化原因，不能只凭现象断言是哪一种。
\item 第二项成立，但要保留\kw{有限数据集}条件。给定一侧 $u\ge0$、另一侧 $u<0$，有限个负类分数的最大值 $a$ 仍小于0；把分界水平从0移到 $a/2$，即可使两侧都严格分开。
\item 第三项错误。感知机收敛定理保证在有限、线性可分训练集上找到分隔参数；它没有以间隔为优化目标。最大间隔要查背面 SVM。
\end{enumerate}
\ans{XOR 的四条线性判别约束相互矛盾，因此一个感知机无法实现。感知机收敛须有可分解；该保证既不是最大间隔保证，也不是测试集泛化保证。}
\qcheck{把候选参数逐点代入最可靠；只画一条看似正确的线不算证明。OR 可由 $w=(1,1)\T,b=-1/2$ 严格分开，可用来对照 XOR。}
\pitfall{“线性分类器”允许 bias，因此边界一般不经过原点。证明时必须照标签写不等号，不能把目标0误写成“分数等于0”。}
\transfer{如果换成更高维特征或含非线性隐藏层，先在新表示空间重新写约束；原始输入空间不可分，并不推出任何模型都不可分。纯线性或仿射层叠加仍只有仿射分数。}

\problem{Q15}{输入空间、权空间与容量计数}{正面第2栏／P02}
\qtext{对带 bias 的二维感知机，说明固定权重与固定样本时分别在什么空间得到超平面。将 OR 四点的分类要求写成权空间约束。设 $m$ 个过原点的超平面位于 $D$ 维权空间且处于一般位置，其最大开区域数为 $R(m,D)=2\sum_{k=0}^{D-1}\binom{m-1}{k}$。计算 $R(4,3)$；解释为什么两个二值输入的16个布尔函数中只有14个为阈值函数，并比较三输入的256个布尔函数、104个阈值函数和128区上界。}
\route{圈出“\kw{input/weight space、duality、regions、Boolean/threshold functions}”，查 P02。这里的对偶是交换“固定样本/固定参数”的几何观点；它与 SVM 的 Lagrangian 对偶不是同一推导。}
\given{输入数 $d=2$，$\tilde x=(x_1,x_2,1)\T$，$\tilde w=(w_1,w_2,b)\T$，参数空间维数 $D=d+1=3$。每个固定样本给一个过权空间原点的平面 $\tilde w\T\tilde x=0$。组合数在 $k>m-1$ 时取0。}
\solve
固定 $\tilde w$ 而改变 $x$，得到输入空间中的边界 $w_1x_1+w_2x_2+b=0$。固定 $x$ 而改变参数，得到权空间中的平面；目标标签选择其哪一侧。OR 的可行参数交集为
\[
 b<0,\quad w_2+b\ge0,\quad w_1+b\ge0,\quad w_1+w_2+b\ge0.
\]
例如 $(1,1,-1/2)$ 位于这个交集中。感知机训练是在权空间寻找一个可行点；输入点本身没有被更新。

对三维中心排列，
\[
 R(4,3)=2\left[\binom30+\binom31+\binom32\right]
 =2(1+3+3)=14.
\]
这里的等号要求一般位置：法向量没有额外退化，使每次切割达到最大数；否则实际区域数至多为该数。四平面都过原点的条件不能删掉，不能把该式当任意位置平面的分区公式。

两个二值输入有 $2^2=4$ 行真值表，每行独立选0或1，所以共有 $2^4=16$ 个函数。权空间每个开区域给一种严格符号组合；这四点的增广平面达到14区，缺少的两种就是 XOR 和 XNOR。

对 $d=3$，输入组合数 $m=8$，含偏置参数维数 $D=4$，故
\[
 R(8,4)=2(1+7+21+35)=128<2^8=256.
\]
课件给实际阈值函数数104。立方体顶点的增广法向量不满足所需的一般位置条件，所以上界128不必达到；正确关系为 $104\le128<256$。
\ans{输入空间边界固定参数、权空间边界固定样本。$R(4,3)=14$；布尔函数总数 $2^{2^d}$，阈值函数只占子集，且 $T_d\le R(2^d,d+1)$。}
\qcheck{维数使用 $D=d+1$，故二维输入不能误代 $R(4,2)=8$。三输入的8是样本组合数，256才是布尔函数数，104是课件列出的实际阈值函数数。}
\pitfall{最大区域数、某组样本的实际区域数和全部布尔函数数是三种量。“更多训练轮次”只能寻找已有可行区域，不能增加该结构能表达的标签模式。}
\transfer{换 $m,D$ 后，组合求和步骤不变；先重新判断是否为中心排列、是否一般位置。输入数增加时参数维数也变，不能把“固定 $D$ 时随 $m$ 为 $D-1$ 次多项式”无条件跨维度套用。}

\problem{Q16}{阈值、bias 与单极／双极编码}{正面第2栏／P01、P02}
\qtext{OR 用 $x_i\in\{0,1\}$ 表示，分数为 $u=x_1+x_2-1$，输出 $z=H(u)$，$H(0)=1$。将输入改为 $s_i=2x_i-1\in\{-1,1\}$，目标改为 $t=2z-1$，写出等价分数与输出。再用双极感知机更新 $\tilde v\leftarrow\tilde v+\eta(t-\hat t)\tilde s$：初参数全0，$\eta=1/2$，先看原始输入 $(0,0)$，完成一次更新。最后判断“双极编码一定有更高准确率”。}
\route{圈出“\kw{unipolar/bipolar、threshold、bias、encoding}”，查 P01/P02。先转换输入坐标和目标编码，再套相应误差；不能只把0替换为−1而保留全部旧参数。}
\given{$x_i=(s_i+1)/2$；$\tilde s=(s_1,s_2,1)\T$，$\tilde v=(v_1,v_2,b_s)\T$；双极输出约定为 $\hat t=2H(u)-1$，故 $u=0$ 时输出 $+1$。}
\solve
代入坐标变换：
\[
 u=\frac{s_1+1}{2}+\frac{s_2+1}{2}-1
 =\frac12s_1+\frac12s_2.
\]
因此等价参数是 $(v_1,v_2,b_s)=(1/2,1/2,0)$，输出为 $\hat t=2H((s_1+s_2)/2)-1$。四个原输入转换成 $(-1,-1),(-1,1),(1,-1),(1,1)$，分数依次为 $-1,0,0,1$，目标为 $-1,+1,+1,+1$，保持相同分类。

训练小问另从\kw{全零参数}出发，并非从上述等价参数出发。第一样本为 $\tilde s=(-1,-1,1)\T$、$t=-1$。旧分数0，预测 $\hat t=+1$，所以
\[
 t-\hat t=-2,\qquad
 \tilde v'=(0,0,0)\T+\frac12(-2)(-1,-1,1)\T=(1,1,-1)\T.
\]
新分数为 $1(-1)+1(-1)-1=-3$，第一样本已正确。

“双极编码一定更准”不成立。编码会改变几何参数化、误差步幅和不同分类在权空间中的比例，可能有代数便利；它不会让 XOR 从不可分变可分，也不能独立保证优化或测试表现。
\ans{等价 OR 分数为 $(s_1+s_2)/2$，双极输出为 $2H(u)-1$。指定的一次更新得到 $(1,1,-1)$；编码本身没有普遍准确率保证。}
\qcheck{输入与输出变换都可逆：$x=(s+1)/2$、$z=(t+1)/2$；逐点回代能检验是否保留边界归属。}
\pitfall{双极误分类时 $t-\hat t=\pm2$。若改用误分类更新 $\eta t\tilde s$，步长常数不同，不能在同一数值学习率下声称更新相同。圆圈内阈值为 $\theta$ 时偏置为 $-\theta$。}
\transfer{对一般 $u=w\T x+b$，输入改为 $s=2x-\mathbf1$ 后有 $v=w/2$、$b_s=b+\frac12\sum_iw_i$。换三输入多数函数时仍先逐坐标代入，再判断阈值。}

\problem{Q17}{ADALINE：连续残差的一次 LMS 更新}{正面第3栏／A01、D02}
\qtext{ADALINE 使用 $a=w\T x+b$，目标 $z=1$，单样本损失 $E=\frac12(z-a)^2$。已知 $x=(1,2)\T$、$w=(0.1,-0.1)\T$、$b=0$，$\eta=0.1$。手推梯度并一次同步更新全部参数，计算新输出与新损失。若最后用 $\operatorname{sign}(a)$ 分类，分类正确后是否应停止所有 LMS 更新？}
\route{圈出“\kw{ADALINE/LMS、linear output、squared error}”，查 A01；D02 区分两种误差。更新依据是目标减\kw{线性输出}，不是目标减硬阈值类别。即使题目最后要求分类，训练仍可最小化连续平方残差。}
\given{$x,w\in\R^2$，$a,z,b,E$ 为标量；$e=z-a$。本题有 $1/2$，所以 $\nabla_w E=-ex$、$\partial E/\partial b=-e$，无额外2因子；线性输出斜率为1。}
\solve
先用旧参数前向：
\[
 a=0.1\times1-0.1\times2=-0.1,\quad e=1-(-0.1)=1.1,
 \quad E=\frac12(1.1)^2=0.605.
\]
用链式法则逐项看，例如 $\partial E/\partial w_2=(a-z)\partial a/\partial w_2=-1.1\times2=-2.2$。因此
\[
 \nabla_w E=(-1.1,-2.2)\T,\qquad \frac{\partial E}{\partial b}=-1.1.
\]
同一旧残差驱动所有更新：
\[
 \begin{aligned}
 w_1'&=0.1-0.1(-1.1)=0.21,\\
 w_2'&=-0.1-0.1(-2.2)=0.12,\\
 b'&=0-0.1(-1.1)=0.11.
 \end{aligned}
\]
全部更新完成后再前向：
\[
 a'=0.21\times1+0.12\times2+0.11=0.56,
 \quad e'=0.44,\quad E'=\frac12(0.44)^2=0.0968.
\]
旧分数为负、新分数为正，类别已从−1变成+1，但 $a'=0.56\ne z=1$，平方损失梯度仍非零。分类判对与残差为0是两个条件。
\ans{新参数 $w'=(0.21,0.12)\T,b'=0.11$，新输出0.56，损失从0.605降至0.0968；分类正确并不意味着 LMS 残差已为0。}
\qcheck{直接回代新参数核对输出；也可对 $E$ 作小扰动检查梯度符号。本样本损失下降不代表其他样本的损失也下降。}
\pitfall{若用硬分类差 $1-(-1)=2$ 更新，就改成了另一套规则。若题目给 $E=(z-a)^2$，梯度应多乘2；若 $a=\kappa(w\T x+b)$，梯度还须乘 $\kappa$。}
\transfer{换样本或目标，仍按“线性输出→连续残差→梯度→同步更新→新输出”执行。换成 batch 时，应在同一旧参数处算全体残差，累加梯度后只更新一次。}

\problem{Q18}{Batch、online 与“最小平方仍做不了 XOR”}{正面第3栏／A01；正面第2栏／P02}
\qtext{(a) 只有可训练 bias，模型 $a=b$，初值 $b=0$；两个目标依次为1、2，$\eta=0.1$。分别用 online 单样本半平方、batch 半平方总和、batch 半平方均值完成一轮更新。(b) 对 XOR，增广设计矩阵与目标为
\[
 X=\begin{bmatrix}0&0&1\\0&1&1\\1&0&1\\1&1&1\end{bmatrix},\quad
 z=\begin{bmatrix}-1\\1\\1\\-1\end{bmatrix}.
\]
用 $E=\frac12\norm{z-X\tilde w}^2$ 推导驻点并求解，解释结果是否意味着 XOR 已被分类学会。}
\route{圈出“\kw{online/batch、sum/mean、least squares、stationary}”，查 A01。第一问比较更新时机及缩放；第二问用线性输出平方目标，不能使用感知机误分类更新。P02 用来检查结构能否实现标签。}
\given{(a) 单样本梯度 $\partial E_p/\partial b=b-z_p$。(b) $X\in\R^{4\times3}$，$\tilde w=(w_1,w_2,b)\T\in\R^3$，$z\in\R^4$；$X\tilde w$ 是四个连续预测。}
\solve
(a) Online 更新两次，第二个残差使用更新后的 bias：
\[
 b_1=0+0.1(1-0)=0.1,\qquad b_2=0.1+0.1(2-0.1)=0.29.
\]
Batch 总损失在旧 $b=0$ 处累加，得 $b'=0+0.1[(1-0)+(2-0)]=0.3$。Batch 均值梯度再除样本数2，得 $b'=0.15$。0.29与0.3来自\kw{残差计算时机}不同；0.3与0.15来自\kw{目标缩放}不同。

(b) 先对参数求导：
\[
 \nabla E=X\T(X\tilde w-z)=0
 \quad\Longrightarrow\quad X\T X\tilde w=X\T z.
\]
逐项内积得到
\[
 X\T X=\begin{bmatrix}2&1&2\\1&2&2\\2&2&4\end{bmatrix},\qquad
 X\T z=\begin{bmatrix}0\\0\\0\end{bmatrix}.
\]
三条方程为 $2w_1+w_2+2b=0$、$w_1+2w_2+2b=0$、$2w_1+2w_2+4b=0$。第一条减第二条得 $w_1=w_2=s$；第一条给 $3s+2b=0$；第三条给 $4s+4b=0$，即 $b=-s$。代回得 $s=0$，所以 $\tilde w=0$。

矩阵 $X$ 满列秩，Hessian $X\T X$ 正定，因此这是唯一全局最小值。预测全为0，残差仍为 $(-1,1,1,-1)\T$，$E=2$；无论把零分数统一归到哪一类，都不可能正确分开 XOR。
\ans{(a) Online 得0.29，batch 总和得0.3，batch 均值得0.15。(b) 最小平方唯一最优为 $w_1=w_2=b=0$；优化已解决该目标，结构仍无法表示 XOR 分类。}
\qcheck{把零向量代入正规方程满足驻点；$X\T z=0$ 表明最优残差与设计矩阵列正交。它不要求每个残差都为0。}
\pitfall{不能写“ADALINE只能batch”：课件展示batch版本，online LMS同样可由单样本平方损失求导。驻点一般不自动是最小值；本题依赖凸平方目标和满列秩作结论。}
\transfer{换数据后重新计算 $X\T X$ 与 $X\T z$；若列不独立，最优系数可能不唯一。换总和/均值时相应调整梯度尺度；固定步长online也不保证每轮总损失单调下降。}

\problem{Q19}{Perceptron V.S. ADALINE：可直接写的比较}{正面第3栏／D02、A01；正面第2栏／P01}
\qtext{用完整句子比较 Perceptron 与 ADALINE，必须涉及共同点、输出/损失/更新依据、训练方式、适用任务和限制。再判断：某题给目标 $z=1$、当前线性分数 $u=0.2$；是否“两个模型已预测正确，所以都不再更新”？假设感知机用 $H(u)$，ADALINE 用半平方损失。}
\route{圈出“\kw{compare/difference/common、hard limiter、linear residual}”，查 D02。简答题先选明确比较维度，再组织句子。看到“学习率/更新权重”不足以识别模型，须继续检查输出和损失。}
\given{两者共同有 $u=w\T x+b$；感知机输出 $\hat z=H(u)$，ADALINE 的训练输出为 $a=u$。目标使用0/1编码，$H(0)=1$；$\eta>0$，若 $x\ne0$，非零残差通常产生非零权重更新。}
\solve
可写成以下完整比较：

\kw{共同点：}两者都是用监督目标调整权重的单层前馈模型，均先形成仿射加权和并可含 bias；若最终用一个阈值分类，原输入中的决策边界都是超平面。

\kw{输出与学习依据：}感知机把加权和送入硬阈值，按目标类别与预测类别的差纠错，已分类正确的样本不触发该规则更新。ADALINE 在训练时保留线性连续输出，最小化平方残差；即使类别正确，只要分数尚未达到目标，梯度仍可非零。

\kw{优化与训练方式：}课件所列感知机规则不靠硬阈值输出误差的普通平滑求导；ADALINE/LMS 由可微平方目标的梯度下降得到。两者均可安排逐样本或整批更新；课件“ADALINE only batch”不能作为一般限制。总和与均值损失需协调学习率。

\kw{任务与限制：}感知机用于线性阈值分类，有限可分训练集有标准收敛保证；ADALINE 可用于线性回归，也可对连续分数阈值化作分类。二者在原始特征上的单一仿射分数都无法正确表达 XOR；ADALINE 的最优平方误差不等于零分类错误。

对指定样本，感知机预测 $H(0.2)=1$，类别差0，故无更新。ADALINE 残差 $e=1-0.2=0.8$，因此 $\Delta w=0.8\eta x$、$\Delta b=0.8\eta$，仍更新。
\ans{共同点是监督、单层前馈和仿射加权和；主要差异是硬类别纠错与线性平方残差梯度。当前感知机不更新，ADALINE 仍按残差0.8更新。自查关键词：\kw{hard limiter / linear output / class error / squared residual / supervised / linear boundary}。}
\qcheck{把两套更新中的“误差”明确写出来：$1-H(0.2)=0$ 与 $1-0.2=0.8$，差异立即可见。}
\pitfall{只说“都是神经网络”不足以比较；说“ADALINE用梯度所以能解所有非线性分类”也错误。不要把 batch 误称为“同时把每个参数逐个改一次后重算残差”。}
\transfer{如果题目换成±1标签，重新注明阈值和误差幅度；比较维度不变。若模型先换了非线性特征，要说明线性边界是在特征空间，而不再断言原输入空间只有直线。}

\problem{Q20}{MLP 前向与 XOR 的两种构造}{正面第3栏／B01；正面第2栏／P02}
\qtext{所有阈值单元使用 $H(u)=1$ 当 $u\ge0$，否则0。网络 A 定义
\[
 h_1=H(x_1-x_2-1/2),\quad h_2=H(-x_1+x_2-1/2),\quad
 y=H(h_1+h_2-1/2).
\]
对四个二值输入计算隐藏状态与输出，说明它为何能实现 XOR。网络 B 只用一个隐藏单元 $h=H(x_1+x_2-3/2)$，输出允许直接连接原输入：$y=H(x_1+x_2-2h-1/2)$。检查 B 的四点输出，并回答：若所有隐藏激活改成线性函数，只加更多层能否解决 XOR？}
\route{圈出“\kw{MLP、hidden layer、forward pass、XOR、skip connection}”，查 B01。这里只算固定参数的推理，无需 BP；BP 是求训练梯度的方法。P02 说明单一超平面为什么不够。}
\given{$x=(x_1,x_2)\T$ 是输入；A 有2个非线性隐藏单元和1个输出单元，两层可训练计算。若计入只传值的输入层，画图会有3层，层数必须说明计法。式中 $-1/2$ 是 bias，不是正阈值。}
\solve
网络 A 按“两个隐藏加权和→两个阈值输出→输出层”计算：
\[
\begin{array}{c|rr|cc|r|c}
x&u_{h_1}&u_{h_2}&h_1&h_2&u_y&y\\\hline
(0,0)&-1/2&-1/2&0&0&-1/2&0\\
(0,1)&-3/2&1/2&0&1&1/2&1\\
(1,0)&1/2&-3/2&1&0&1/2&1\\
(1,1)&-1/2&-1/2&0&0&-1/2&0
\end{array}
\]
$h_1$ 检测“$x_1$ 比 $x_2$ 大”，$h_2$ 检测反方向，输出对两种“不同”作 OR。原空间需要两个区域；隐藏空间中正类为01/10，负类为00，用 $h_1+h_2\ge1/2$ 一次分开。

网络 B 的隐藏单元是 AND，四点 $h=0,0,0,1$。输出分数依次为 $-1/2,1/2,1/2,-1/2$，仍得 XOR。它只有两个计算单元，但有\kw{输入直接到输出的跨层连接}，不能把它误画成普通无跨层连接的一个隐藏节点网络。

若隐藏层仅作仿射变换，$h=Ax+c$，输出激活前分数 $v\T h+d=v\T Ax+v\T c+d$ 仍是仿射函数。任意多层仿射变换复合仍仿射；末尾一个阈值仍只有一个超平面，不能突破 XOR 限制。
\ans{A 与 B 都输出 $0,1,1,0$。A 通过非线性隐藏表示把 XOR 变为可分；B 还依赖跨层连接。增加纯线性/仿射层不能解决 XOR。}
\qcheck{四点真值表是本题完整检验；不能只算输入 $(1,0)$ 就宣称网络实现 XOR。A 的隐藏输出维数2与输入维数2相同，非线性改变表示不一定升维。}
\pitfall{普通 BP 不能对硬阈值按 sigmoid 导数求导；本题是固定逻辑网络的前向计算。隐藏层没有题目直接给出的目标，不表示它的活动不能被计算或观察。}
\transfer{换权重或阈值，前向顺序不变，但需重新逐点检查。换为 sigmoid 网络时，输出通常为连续近似值，应分清实际激活与阈值化后的类别，不能把0.9直接写成精确布尔1。}

\problem{Q21}{完整一次 BP：四个参数全部同步更新}{正面第4栏／B02；激活见正面第3栏／B01}
\qtext{一个输入、一个隐藏节点和一个输出节点：
\[
 h=\sig(q),\quad q=wx+b_h,\qquad
 o=\sig(r),\quad r=vh+b_o,\qquad \sig(u)=\frac1{1+e^{-u}}.
\]
取 $x=z=1$，初值 $(w,b_h,v,b_o)=(0,0,2,-1)$，$E=\frac12(z-o)^2$，$\eta=1$，无 momentum。按课件约定 $\delta=-\partial E/\partial(\text{激活前输入})$，完成前向、反向、四个梯度及一次更新。}
\route{圈出“\kw{MLP、sigmoid、half squared error、backpropagation、one update}”，查 B02。B01 提供局部导数。与 LMS 不同，这里需把输出残差经输出激活、连接权及隐藏激活依次传回；直接把残差乘输入会漏因子。}
\given{$x,z,w,b_h,v,b_o,h,o$ 都是标量。$w$ 是输入到隐藏权，$v$ 是隐藏到输出权。本题统一负导数 $\delta_o=-\partial E/\partial r$、$\delta_h=-\partial E/\partial q$，所以参数更新为加 $\eta\delta\times$前层活动。}
\solve
\kw{1. 旧参数前向并记录活动。}
\[
 q=0\times1+0=0,\quad h=\frac12,\qquad
 r=2\times\frac12-1=0,\quad o=\frac12,\quad E=\frac18.
\]
\kw{2. 输出层链式求导。}三个因子为残差方向、局部激活导数、前层活动：
\[
 \delta_o=(z-o)o(1-o)=\frac12\cdot\frac14=\frac18,
\]
\[
 \frac{\partial E}{\partial v}=-\delta_o h=-\frac1{16},\qquad
 \frac{\partial E}{\partial b_o}=-\delta_o=-\frac18.
\]
\kw{3. 隐藏层接收下游贡献。}此时仍用旧连接权 $v=2$：
\[
 \delta_h=h(1-h)v\delta_o=\frac14\cdot2\cdot\frac18=\frac1{16},
\]
\[
 \frac{\partial E}{\partial w}=-\delta_h x=-\frac1{16},\qquad
 \frac{\partial E}{\partial b_h}=-\delta_h=-\frac1{16}.
\]
\kw{4. 梯度全部算完后统一更新。}
\[
 \boxed{(w',b_h',v',b_o')=\left(\frac1{16},\frac1{16},\frac{33}{16},-\frac78\right).}
\]
例如 $v'=2-(-1/16)=33/16$。新参数用于下一次前向，而不用于重写本次隐藏梯度。
\ans{旧 $h=o=1/2$、$E=1/8$；$\delta_o=1/8,\delta_h=1/16$。四个梯度依 $(w,b_h,v,b_o)$ 顺序为 $(-1/16,-1/16,-1/16,-1/8)$，更新参数如框中所示。}
\qcheck{新 $q'=1/8$，故 $h'=\sig(1/8)$；$o'=\sig((33/16)\sig(1/8)-7/8)\approx0.5549$，$E'\approx0.0990<0.125$。小扰动中心差分也可核对四个旧梯度。}
\pitfall{若先更新 $v$ 再算 $\delta_h$，就混用两套参数。若把 $\delta$ 定义为正导数，公式符号会整体变化；不能正导数配“加 $\eta\delta$”。bias 的梯度无需再乘 $h$ 或 $x$。}
\transfer{多输出时，隐藏节点把所有下游 $w\delta$ 相加；batch 时先在共同旧参数处累加各样本梯度。换激活应重算局部导数；换成不含 $1/2$ 的平方损失，全部梯度多乘2。}

\problem{Q22}{隐藏层的分支求和与 momentum}{正面第4栏／B02、B03}
\qtext{某 logistic 隐藏节点当前活动 $h=0.8$，有两个下游连接，旧权分别 $v_1=0.5,v_2=-0.3$，下游负导数 deltas 为 $\delta_1=0.1,\delta_2=-0.2$。该节点一个输入为 $x=-2$，对应旧权 $w=0.4$。先求隐藏 $\delta_h$ 与 $\partial E/\partial w$。再用 $\eta=0.05$、动量系数 $\mu=0.8$、上一步权重增量 $\Delta w_{\rm prev}=0.01$，求新增量与新权重。解释动量是否保证找到全局最小。}
\route{圈出“\kw{hidden、two downstream paths、momentum、previous increment}”，查 B02/B03。隐藏节点没有直接目标，不能自行编一个 $z-h$；必须汇总所有下游链式贡献。动量还需读取上一步\kw{增量}。}
\given{$\delta_j=-\partial E/\partial u_j$；$\delta_h=h(1-h)\sum_jv_j\delta_j$，$\partial E/\partial w=-\delta_h x$。动量规则为 $\Delta w=-\eta\partial E/\partial w+\mu\Delta w_{\rm prev}$，$w'=w+\Delta w$。}
\solve
先算两条分支的加权和：
\[
 v_1\delta_1+v_2\delta_2=0.5(0.1)+(-0.3)(-0.2)=0.11.
\]
隐藏局部导数为 $h(1-h)=0.8\times0.2=0.16$，所以
\[
 \delta_h=0.16\times0.11=0.0176,\qquad
 \frac{\partial E}{\partial w}=-(0.0176)(-2)=0.0352.
\]
虽然隐藏 delta 为正，这条权重的\kw{负输入}使梯度为正；无动量梯度下降会减小该权重。
\[
 \Delta w=-0.05(0.0352)+0.8(0.01)=-0.00176+0.008=0.00624,
 \qquad w'=0.40624.
\]
本次增量反而为正，因为前次方向的累积项更大。动量可平滑方向变化、减轻狭长误差谷中的振荡，但不是全局搜索证明；它也不保证每一步目标下降。
\ans{$\delta_h=0.0176$，梯度0.0352，新增量0.00624，新权重0.40624。动量保留历史运动方向，不保证全局最优。}
\qcheck{删去动量即得到 $w'=0.39824$；与原结果比较可检查两项符号。隐藏分支若只保留第一项，会误得 $\delta_h=0.008$，明显漏掉第二条有效路径。}
\pitfall{不能用 $\mu w$ 代替 $\mu\Delta w_{\rm prev}$。同一步的分支权应为旧参数；负连接权与负 delta 的乘积为正，不可凭箭头符号直觉跳过乘法。}
\transfer{增加下游节点只需扩大求和；换为 tanh 隐藏节点须将局部导数改成 $1-h^2$。若题目给的“动量变量”是速度而非参数增量，先把其递推和符号与本规则对应，避免混用定义。}

\problem{Q23}{激活函数纠错：少一个尺度就改了梯度}{正面第3栏／B01；反传见正面第4栏／B02}
\qtext{判断并改正 $g(u)=(1-e^{-u})/(1+e^{-u})=\tanh(u)$。推导 $g'(u)$ 并用输出 $g(u)=0.6$ 计算局部导数。另有斜率为2的 logistic $a(u)=1/(1+e^{-2u})$，求 $a'(0)$。最后解释：为何硬阈值不能直接沿用这些导数，sigmoid 的饱和会带来什么训练问题？}
\route{圈出“\kw{sigmoid/tanh、derivative、slope、correct the formula}”，查 B01。先看\kw{完整指数和尺度}，不能因为两条曲线都为S形就套同一个导数；B02 中该局部导数会乘入整条链式梯度。}
\given{$\sig(u)=(1+e^{-u})^{-1}$，$\sig'(u)=\sig(u)(1-\sig(u))$；$\tanh(u)=(1-e^{-2u})/(1+e^{-2u})$。这里 $u$ 为激活前输入，$g,a$ 为活动值，均为标量。}
\solve
将 $u/2$ 代入标准 tanh 定义：
\[
 g(u)=\tanh(u/2)=2\sig(u)-1.
\]
由链式法则或商法则可得
\[
 g'(u)=\frac12[1-\tanh^2(u/2)]
 =\frac12(1-g(u)^2).
\]
所以当 $g=0.6$ 时，$g'=\frac12(1-0.36)=0.32$。若误用标准 $\tanh(u)$ 导数，会得到0.64，为正确值的两倍。

斜率2的 logistic 是 $a(u)=\sig(2u)$，故
\[
 a'(u)=2a(u)(1-a(u)),\qquad a(0)=\frac12,\qquad a'(0)=\frac12.
\]
普通 $\sig(u)$ 在0处斜率为1/4，标准 $\tanh(u)$ 在0处为1，$\tanh(u/2)$ 在0处为1/2，可作为三个公式的快速尺度检查。

硬阈值在非边界处导数为0、跳变处不可按普通微分求导，不能直接给它套 $a(1-a)$。Logistic/tanh 在两端饱和，导数接近0，链式传播时小导数连乘可使早层梯度很小，导致学习缓慢。
\ans{原等式应改为 $g(u)=\tanh(u/2)$，导数 $\frac12(1-g^2)$；指定活动处导数0.32；斜率2的 logistic 在0处导数0.5。}
\qcheck{在 $u=0$ 附近取很小的正负扰动，用函数值差商检查斜率；标准 tanh 与本题 $g$ 输出零点相同，但斜率不同。}
\pitfall{饱和区导数\kw{接近}0，不表示 sigmoid 处处导数为0。分段线性压缩函数在折点处也需另定导数，不能把“连续”直接当“处处可微”。}
\transfer{一般 $\sig(cu)$ 多乘 $c$；一般 $\tanh(cu)$ 的导数为 $c(1-a^2)$。换激活只改变局部导数因子，前向记录、反向链式和同步更新的框架保持不变。}

\problem{Q24}{表示能力、优化成功与泛化是三件事}{正面第4栏／B03；模型结构见正面第3栏／B01}
\qtext{有人说：“MLP 是通用逼近器，Kolmogorov 也说明多变量函数可分解，所以固定一个很小的 sigmoid 网络、持续运行 BP，一定会找到全局最优并对新数据正确。”逐项回应。再看真实规律 $f(x)=x$，训练数据只有 $(0,0),(1,1)$；候选为 $g_1(x)=x$ 与 $g_2(x)=x+4x(1-x)$。比较训练误差及 $x=1/2$ 处预测。若训练损失持续下降而验证损失先降后升，应怎样使用验证集与测试集？}
\route{圈出“\kw{universal approximation、exist/find/generalize、overfitting、validation}”，查 B03。这里不是再算一次梯度，而是检查结论声称解决了哪一个问题：存在参数、算法找到参数、未见数据有效。}
\given{表示能力讨论模型家族及容量；优化讨论选定目标、算法和初值；泛化讨论训练与模型选择之外的数据。算术小问统一用单点半平方损失 $\frac12(f-g)^2$，输入域为 $[0,1]$。}
\solve
\begin{enumerate}
\item \kw{表示：}通用逼近结论在适当激活、容量和紧域上的连续目标等条件下，允许选择足够的隐藏单元，使近似误差任意小。它不是固定某个很小网络对任意函数精确相等的结论。Kolmogorov 的存在性分解允许合适的一元函数，不能未经证明把它们全部换成一个指定的小 sigmoid 网络。
\item \kw{优化：}参数存在不代表 BP 从任何初值都能找到。一般多层训练目标非凸，可能存在局部极小、鞍点或平坦区域；驻点也不自动是全局最优。Momentum 能改善某些轨迹，但没有解除上述限制。
\item \kw{泛化：}有限训练点不能唯一确定点间行为。低训练误差没有独自证明对新样本的表现，通用逼近存在性也没有提供这种保证。
\end{enumerate}
本例中，两候选在0、1处都与真实目标相同，训练损失都为0。但在 $x=1/2$，
\[
 f=\frac12,\quad g_1=\frac12,\quad g_2=\frac12+4\cdot\frac12\cdot\frac12=\frac32.
\]
故该点 $g_1$ 损失0，$g_2$ 损失 $\frac12(1)^2=1/2$。这是零训练误差仍可泛化较差的明确例子。

训练/验证损失分离提示过拟合或数据划分问题。用验证集选择训练轮次和模型设置，保留验证表现较好的参数；可比较容量、正则化等设置。测试集应保留为最终评价，不反复据测试结果调参后仍称其为独立检验。
\ans{存在性、优化、泛化分别需要不同证据，三者不能互相替代。本例两模型训练误差均为0，但在 $1/2$ 处 $g_2$ 预测1.5、半平方损失0.5；继续降低训练误差不能保证恢复泛化。关键词：\kw{capacity / nonconvex / validation / overfitting}。}
\qcheck{先用训练点验证“两者都拟合”，再用未参与拟合的点验证差异；不要只凭曲线弯曲与否定义过拟合。}
\pitfall{“网络大”并非过拟合的充分定义；“总损失驻点”也不是“每个样本残差为0”。不能从盲人摸象插画推出课件没有给出的数学定理。}
\transfer{若改成不连续目标或换域，须重新检查逼近定理条件。模型选择程序仍应区分训练、验证、测试；测试好也只支持已进行的评价范围，不能保证任意未来分布。}

\problem{Q25}{RBF 特征值、中心和宽度}{正面第5栏／R01}
\qtext{Gaussian RBF 定义为 $\phi_j(x)=\exp[-\norm{x-c_j}^2/\sigma_j^2]$。取 $c_1=(0,0)\T,c_2=(1,1)\T$，两个宽度都为1。计算 $x=(1,0)\T$ 与 $x=(0,0)\T$ 的特征向量。若同样距离 $r=1$ 的一个基函数宽度由1改为2，响应如何改变？若题目改用 $\exp[-r^2/(2\sigma^2)]$，$r=\sigma=1$ 时又是多少？说明“径向”与“Gaussian 局部响应”的区别。}
\route{圈出“\kw{RBF、center、distance、width、Gaussian}”，查 R01。先核对分母约定，再算平方距离和指数。这里是在计算显式特征；仅见 Gaussian 并不能认定题目正在训练 SVM。}
\given{$x,c_j\in\R^2$，$r_j=\norm{x-c_j}\ge0$，$\sigma_j>0$；$\phi(x)=(\phi_1,\phi_2)\T\in\R^2$。输出权重尚未给出，因此本题只算特征，不能自行把特征和称为预测。}
\solve
对 $x=(1,0)\T$，两个平方距离均为1：
\[
 \norm{x-c_1}^2=1^2+0^2=1,\qquad
 \norm{x-c_2}^2=0^2+(-1)^2=1,
 \quad\phi(x)=(e^{-1},e^{-1})\T\approx(0.3679,0.3679)\T.
\]
对 $x=(0,0)\T$，到两个中心的平方距离为0、2，所以
\[
 \phi(x)=(1,e^{-2})\T\approx(1,0.1353)\T.
\]
固定 $r=1$ 时，宽度由1变2使响应由 $e^{-1}\approx0.3679$ 变为 $e^{-1/4}\approx0.7788$。宽度越大，离中心同样远的位置仍可获得较大响应；不是峰值变高，中心处峰值仍为1。

若分母给 $2\sigma^2$，则响应为 $e^{-1/2}\approx0.6065$，应照题使用；两个宽度参数的数值含义不同。

“径向”表示只依赖距离 $r$；同中心、同半径的位置有相同响应。课件还列 $r^p$（$p=1,3,\ldots$）、薄板 $r^2\log r$、多二次 $\sqrt{r^2+\sigma^2}$ 及其倒数。它们也径向，但不都像 Gaussian 那样远离中心趋于0；薄板式在 $r=0$ 按极限取0。
\ans{两个输入的特征分别为 $(e^{-1},e^{-1})\T$ 与 $(1,e^{-2})\T$。宽度1变2，指定响应由0.3679变0.7788；分母 $2\sigma^2$ 的响应为0.6065。Gaussian 的局部钟形是特定基函数的性质。}
\qcheck{中心处平方距离0、响应1；有限距离时 Gaussian 响应应在 $(0,1]$。两个中心等距时两响应相等。}
\pitfall{题目已经给 $\norm{x-c}^2$，不要再平方成距离四次方；$r$ 是距离，$r^2$ 才是平方距离。宽度参数要正；不要把分母里的 $\sigma^2$ 与 $2\sigma^2$ 混用。}
\transfer{换中心会改变每点到中心的距离；换宽度只改指数尺度。中心与宽度固定后可查 R01 解线性输出；若同时优化中心/宽度，整体训练就不再只是线性最小二乘。}

\problem{Q26}{固定 RBF 特征后求输出权重}{正面第5栏／R01；矩阵消元见正面第1栏／M03}
\qtext{一维 RBF 网络有中心 $c_1=0,c_2=1$，$\sigma^2=1/\log2$，$\phi_j(x)=e^{-(x-c_j)^2/\sigma^2}$。输出 $f(x)=a_1\phi_1(x)+a_2\phi_2(x)$，\kw{不另加 bias}。要求在 $(x,z)=(0,1),(1,0)$ 两点精确插值。构建特征矩阵、消元求 $a_1,a_2$，预测 $x=1/2$。解释为何若额外加入自由 bias，不能继续声称两条方程唯一决定三个未知数。}
\route{圈出“\kw{fixed centers/width、linear output、interpolation、solve weights}”，查 R01。本题未知量是输出系数，特征已固定，适合线性方程求解；无需对中心做 BP，也不需要 SVM 的间隔约束。}
\given{两样本 $N=2$、两基函数 $h=2$，$\Phi\in\R^{2\times2}$，$a=(a_1,a_2)\T\in\R^2$，$z=(1,0)\T$。矩阵第 $i$ 行是第 $i$ 个样本的特征，求解 $\Phi a=z$。}
\solve
距离平方为0时响应1，为1时响应 $e^{-\log2}=1/2$，故
\[
 \Phi=\begin{bmatrix}1&1/2\\1/2&1\end{bmatrix},\qquad
 \begin{bmatrix}1&1/2&\vline&1\\1/2&1&\vline&0\end{bmatrix}.
\]
作行变换 $R_2\leftarrow R_2-\frac12R_1$：
\[
 \begin{bmatrix}1&1/2&\vline&1\\0&3/4&\vline&-1/2\end{bmatrix}.
\]
回代得 $a_2=-2/3$，再由第一行 $a_1=1-\frac12(-2/3)=4/3$。

预测点到两中心的平方距离均为 $1/4$，所以
\[
 \phi_1(1/2)=\phi_2(1/2)=e^{-(\log2)/4}=2^{-1/4},\qquad
 f(1/2)=\left(\frac43-\frac23\right)2^{-1/4}\approx0.5606.
\]
虽然每个 Gaussian 响应非负，输出权重可以为负。

若另加自由 bias $b$，矩阵要增广一列1，未知变为 $(a_1,a_2,b)$。两行三列不可能满列秩；存在插值解时通常不唯一，需要更多独立数据或另外给定选择/正则化原则。不能把原二阶矩阵的解直接当新模型唯一解。
\ans{$a=(4/3,-2/3)\T$，$f(1/2)=\frac23\,2^{-1/4}\approx0.5606$。固定特征后，对输出权重是线性求解，对原输入仍是非线性函数。}
\qcheck{$f(0)=4/3-(2/3)(1/2)=1$，$f(1)=(4/3)(1/2)-2/3=0$；$\det\Phi=3/4\ne0$，本题插值系数唯一。}
\pitfall{矩阵行是样本、列是特征；中心数量与样本数量概念不同。只因本题方阵可逆，不能推出所有 RBF 训练都能直接求逆；一般近似问题是 $\min_a\norm{z-\Phi a}^2$。}
\transfer{换目标仅改右端；换中心/宽度须重算特征矩阵。若数据更多，通常解最小二乘而非强制插值；若特征列相关，正规方程可能非唯一，应按题目指定的伪逆或正则化规则处理。}

\problem{Q27}{Cover 思想：XOR 换表示后怎样可分？}{正面第5栏／R01；线性不可分见正面第2栏／P02}
\qtext{XOR 标签为00/11负类，01/10正类。按新课件中心顺序固定 $c_1=(1,1)\T,c_2=(0,0)\T$，$\phi_j(x)=e^{-\norm{x-c_j}^2}$。重算四点特征；检验 $g(x)=0.9-\phi_1(x)-\phi_2(x)$ 能否分开两类。写出带 bias 的 $\phi$ 可分条件，并纠正“任何升维映射都会可分、会提高泛化”。}
\route{圈出“\kw{Cover、$\phi$-separable、XOR、feature space}”，查 R01。先计算映射再检查线性分数；原空间 XOR 不可分不妨碍在合适新特征中可分。此题中心顺序与 Q25 不同，不能直接抄上一题列顺序。}
\given{$x,c_j\in\R^2$，$\phi(x)\in\R^2$；输出分数的权重为 $a=(-1,-1)\T$，bias为0.9。需要齐次形式时用 $\tilde\phi=(\phi_1,\phi_2,1)\T$、$\tilde a=(-1,-1,0.9)\T$。}
\solve
令 $q=e^{-1}\approx0.367879$、$r=e^{-2}\approx0.135335$。先按两个中心顺序计算平方距离：
\[
\begin{array}{c|cc|cc|c}
x&\norm{x-c_1}^2&\norm{x-c_2}^2&\phi_1&\phi_2&g\\\hline
(0,0)&2&0&r&1&0.9-(1+r)\\
(0,1)&1&1&q&q&0.9-2q\\
(1,0)&1&1&q&q&0.9-2q\\
(1,1)&0&2&1&r&0.9-(1+r)
\end{array}
\]
负类分数约为 $-0.235335$，正类分数约为 $0.164241$，均严格符合符号。因此在新坐标中，直线 $\phi_1+\phi_2=0.9$ 分开了两类。

本例两个不同的正类输入01与10映到同一个特征点，但标签相同，并不妨碍分类。新映射还是二维到二维；起作用的是\kw{非线性改变表示}，不只是增加坐标数量。原输入中的边界变成两个指数函数之和等于0.9，一般不再是一条直线。

有 bias 的 $\phi$ 可分性指存在 $a,b$ 使正类 $a\T\phi(x)+b>0$、负类 $a\T\phi(x)+b<0$；用增广常数坐标可写成齐次形式。课件以 Cover 为标题主要给出这个定义及表示思想，不应声称已展示了完整定理证明。

反例：把任意输入映为零向量，再增加任意多个零坐标，异类样本仍完全重合，任何确定的线性输出都无法区分。即使训练集变可分，也没有因此自动获得测试集保证。
\ans{指定映射后 XOR 严格可分；00/11分数约−0.2353，01/10约0.1642。适当非线性表示可以提高可分性，任意映射或任意升维没有此保证，泛化还需独立评价。}
\qcheck{与课件p146一致时，00必须映为 $(e^{-2},1)$；p147原表及图交换了两特征列。中心、列、图轴与权重必须同步，不能拼接不一致的定义。}
\pitfall{线性可分须说清在哪个空间。不要把本题0.9分类线称为已求得 SVM 最大间隔线；本题只验证符号，不做间隔优化。}
\transfer{交换两中心时，应同时交换特征列和对应输出权重；本题权重相同，预测恰好不变。若权重不同或宽度改变，须重新计算分数。分类正确也不等于精确插值到0/1目标。}

\problem{Q28}{FLN 与 ELM：特征从哪里来，求哪些参数？}{正面第5栏／R02；比较见 D04}
\qtext{(a) 二值输入的 FLN 使用固定特征 $(x_1,x_2,x_1x_2)$，证明输出 $f=x_1+x_2-2x_1x_2$ 精确等于 XOR，说明训练的是哪些量。(b) 基本 ELM 的一次初始化已给定并固定隐藏特征 $\phi_1(x)=\sig((\log3)x)$、$\phi_2(x)=\sig(- (\log3)x)$，$\sig(u)=(1+e^{-u})^{-1}$。不加输出 bias，输出 $f=\beta_1\phi_1+\beta_2\phi_2$；样本为 $(x,z)=(0,0),(1,1)$。构建 $H\beta=z$，求输出权并预测 $x=-1$。解释为什么 ELM 不是“没有隐藏层”。}
\route{圈出“\kw{functional link、fixed expansion、random fixed hidden、output weights}”，查 R02。FLN 的函数项预先指定；基本 ELM 隐藏权和 bias 初始化后固定。此时只有输出系数需用标签拟合，不能把隐藏参数也当本次线性方程的未知数。}
\given{(a) 原输入 $x\in\R^2$，展开特征 $\phi\in\R^3$，系数为 $(1,1,-2)$，bias为0。(b) $N=h=2$，$H_{ij}=\phi_j(x_i)$，$H\in\R^{2\times2}$，$\beta\in\R^2$，$z=(0,1)\T$；隐藏斜率 $\pm\log3$ 与隐藏 bias 0 已固定。}
\solve
(a) 逐点代入：$f(0,0)=0$、$f(0,1)=1$、$f(1,0)=1$、$f(1,1)=2-2=0$。乘积特征记录两输入共同出现的作用。固定展开后，输出对\kw{系数}线性，对\kw{原输入}可以非线性；学习时拟合输出系数及题目允许的 bias，不必把函数项本身也当可训练隐藏层。

(b) $\sig(0)=1/2$，$\sig(\log3)=3/4$、$\sig(-\log3)=1/4$，所以
\[
 H=\begin{bmatrix}1/2&1/2\\3/4&1/4\end{bmatrix},\qquad
 H\beta=\begin{bmatrix}0\\1\end{bmatrix}.
\]
将两式分别乘2、4，得 $\beta_1+\beta_2=0$、$3\beta_1+\beta_2=4$；相减得 $2\beta_1=4$，所以 $\beta_1=2,\beta_2=-2$。在 $x=-1$，特征为 $(1/4,3/4)$，故 $f(-1)=2/4-6/4=-1$。

基本 ELM 含一个隐藏层，通常随机产生并固定隐藏输入权和 bias，再用最小二乘/伪逆等方法求输出权。课件“one-layer”不能解读成没有隐藏节点；\kw{一层可拟合权重}与\kw{一个隐藏层}不是相同的计层口径。固定 $H$ 的线性平方子问题没有坏的局部极小，但特征质量、秩、节点数和泛化仍需判断。
\ans{FLN 的固定乘积特征使 XOR 对系数线性可表示；ELM 本题 $H$ 如上，$\beta=(2,-2)\T$，$f(-1)=-1$。两者主要区别在特征来源，不在“一个非线性、另一个线性”的二分。}
\qcheck{$H(2,-2)\T=(0,1)\T$；$\det H=-1/4\ne0$，本题输出系数唯一。预测−1只是该模型在给定初始化和数据下的结果，不能把它当已获保证的真实规律。}
\pitfall{“随机”不表示推理时重新抽隐藏权。若隐藏权在输出训练期间也持续优化，固定特征线性求解的前提已改变；若 $H$ 非方阵或非满秩，不可无条件写 $H^{-1}$。}
\transfer{换标签仅改右端；换初始化、激活或节点数须重建 $H$。一般可写 $\hat\beta=H^+z$，含额外输出 bias 时须在 $H$ 加一列1并重新检查秩；该线性求解不自动替你选择好的特征。}

\problem{Q29}{MLP V.S. RBF：模型与训练方法分开答}{正面第5栏／D03、R01；正面第4栏／B02}
\qtext{比较 MLP 与 RBF 网络的共同点、隐藏特征、输出、训练方式、任务与限制。有人写“BP、MLP、RBF 是三种并列网络；MLP 一定慢而且不能回归，RBF 一定只做无监督聚类”，请纠正。某任务已给定合适中心和宽度，需要快速拟合连续输出，应怎样说明选择理由及保留条件？}
\route{圈出“\kw{MLP versus RBF、common/difference、choose model、BP}”，查 D03。按架构和参数训练角色比较；不要只列“都是AI”，也不要把求导方法与模型名称排成同类。}
\given{典型 MLP 隐藏活动为 $a_j=s(v_j\T x+b_j)$；Gaussian RBF 隐藏活动为 $\phi_j=e^{-\norm{x-c_j}^2/\sigma_j^2}$。比较的“局部性”针对这里的 Gaussian RBF，不是所有径向基都衰减。输出形式和损失应按任务明确。}
\solve
\kw{共同点：}两者都是前馈参数模型，都能先产生非线性中间特征，再形成预测；适当设计输出和损失后，均可用于分类或回归。模型具有表示能力，不等于已优化成功或已泛化。

\kw{隐藏表示差异：}典型 MLP 节点先作仿射投影再激活，方向、偏置和多层组合可由训练学习；常见 sigmoid 单元对整个投影方向产生响应。Gaussian RBF 节点按到中心的距离局部响应，中心决定位置，宽度决定作用范围；输出是径向特征的线性组合。

\kw{训练差异：}MLP 常用 BP 计算所有可训练层的梯度，再配合梯度下降等优化。RBF 常先选中心/宽度，再固定特征求线性输出权；中心选择可以采用无监督步骤，但输出通常使用监督目标，中心也可以另外监督优化。固定特征时的输出平方问题是凸的；同时训练中心/宽度时，不能声称全问题仍只需一次线性求解。

\kw{任务与限制：}MLP 适合学习层次或组合表示，但训练非凸，架构、初始化与正则化影响结果；RBF 的局部表示直观，固定特征后输出求解直接，但需要足够且合适的中心与宽度，矩阵条件和样本覆盖影响效果。无法仅凭名称保证谁始终更快、更准。

纠错：BP 是链式计算梯度的算法，不是与 MLP/RBF 同层级的网络架构。MLP 可以回归，RBF 网络也并非只有无监督聚类。对所给任务，可选\kw{固定中心和宽度的 RBF+线性输出平方拟合}，理由是隐藏特征已定，输出可由线性最小二乘高效拟合；仍须检查特征覆盖、矩阵秩和独立验证表现。
\ans{两者都可非线性前馈预测；典型 MLP 学习仿射激活的组合，Gaussian RBF 使用到中心的局部距离特征。BP 是梯度方法；RBF 固定特征后可直接拟合线性输出，但不保证所有中心设置都好。关键词：\kw{affine activation / radial distance / local response / fixed features / gradient}。}
\qcheck{比较答案同时出现共同点、同一维度的差异、任务和限制；仅写“一个全局，一个局部”会省掉训练条件与应用。}
\pitfall{前馈指连接方向无环，不表示训练时不能从输出向前层反传导数。误差反传不会把 MLP 变成循环网络。}
\transfer{若换成不同 MLP 输出激活或非 Gaussian 径向函数，须重述相应性质；共同的“特征→输出”结构及表示/优化/泛化区分仍可使用。训练速度应结合数据规模、特征数和求解器讨论。}

\problem{Q30}{RBF V.S. FLN V.S. ELM：共同线性输出的条件}{正面第5栏／D04、R01、R02}
\qtext{比较 RBF 网络、FLN 与基本 ELM 的共同结构、特征来源、固定与训练参数、适用任务及限制。判断：(i) 固定 RBF 中心/宽度后求输出权；(ii) 训练中持续移动 RBF 中心；(iii) 固定 FLN 乘积展开后求输出权；(iv) ELM 每次预测都重新随机隐藏权。哪些保留“固定特征线性输出求解”的前提？解释“快速训练、消除局部极小”能说明什么，不能说明什么。}
\route{圈出“\kw{RBF/FLN/ELM、feature source、which parameters、local minima}”，查 D04。模型比较要追踪特征矩阵是否随未知量改变；只看都含非线性函数不足以判断优化是否为线性平方问题。}
\given{共同写法为 $f(x)=a\T\phi(x)+b$；样本排成 $\Phi_{ij}=\phi_j(x_i)$。当 $\Phi$ 固定、目标为平方残差时，输出参数可作线性最小二乘；若含 bias，增广一列1。}
\solve
\kw{共同点：}三者都可显式计算非线性特征，再对输出系数作线性拟合；按输出和目标设计，可用于分类或回归。对原输入非线性与对输出参数线性可以同时成立。

\kw{RBF：}特征依赖到中心的距离，Gaussian 还需正宽度；中心、宽度、基函数与个数需先选或训练。固定这些特征后拟合输出权和可选 bias；其限制包括中心覆盖、宽度和矩阵条件。

\kw{FLN：}特征是原输入加预设函数项，如乘积、平方；展开函数固定后学习输出系数。它省去另行训练隐藏表示，但函数类型与展开规模需要选择，预设项可能不能充分表达任务规律。

\kw{基本 ELM：}隐藏输入权和 bias 随机初始化后固定，由隐藏响应构成 $H$，主要拟合输出权；隐藏节点数和激活仍要选择。固定的一次随机表示可能质量较差或秩不足，随机并不自动等于可分或泛化好。

逐项判断：(i) 保留；(ii) 不保留，因为移动中心改变每个距离和 $\Phi$；(iii) 保留；(iv) 错误，预测应使用训练时固定的隐藏参数，否则已学习的输出系数对应的特征坐标改变。

固定特征且采用线性输出平方损失时，目标为凸二次函数；其驻点是全局最小，满列秩时输出系数唯一，否则可有多个同等最小值。“无坏局部极小”只适用于这个子问题，不替代特征选择，不保证整个中心/宽度训练凸，也不保证最小训练误差等于最佳泛化。
\ans{三者共同是“显式特征→线性输出”；RBF 按距离，FLN 用预设函数，ELM 用随机固定隐藏响应。(i)(iii) 保留固定特征前提；(ii) 改变特征；(iv) 破坏训练/推理一致性。关键词：\kw{centers/widths / prescribed functions / random fixed hidden / linear least squares}。}
\qcheck{检查未知量是否只出现在最后的 $a,b$；若未知中心、宽度或隐藏权出现在 $\phi$ 内，就不能把整题当固定矩阵求一次线性输出。}
\pitfall{“对参数线性”不等于“对输入线性”。固定特征的平方问题即使凸，也可能不唯一；输出无坏局部极小不等于整个模型没有选择困难。}
\transfer{换数据时保持特征机制定义；换中心、展开项或 ELM 初始化必须重建矩阵并重求输出权。若目标换成其他损失或输出再加非线性，须重新核对是否仍可用线性最小二乘。}
